\section{Alignement final DevSecOps}

\subsection{Position de reference}

Le scenario officiel de soutenance de \textit{SecureRAG Hub} est le mode
\texttt{demo}. Jenkins reste l'autorite CI/CD officielle du projet, tandis que
GitHub Actions est conserve comme historique technique. Cette decision evite de
presenter deux chaines concurrentes et clarifie le perimetre de demonstration.

\subsection{Elements validates localement}

Les elements suivants sont consideres comme validates dans l'etat courant du
projet :

\begin{itemize}
  \item execution des tests applicatifs Laravel ;
  \item rendu Kustomize des overlays \texttt{dev} et \texttt{demo} ;
  \item controles statiques Kubernetes : cleartext scope, ressources et manifests Kyverno ;
  \item audit des dependances Composer/npm lorsque les outils sont disponibles ;
  \item Jenkins local Docker ;
  \item campagne CD \texttt{demo} en mode \texttt{dry-run} ;
  \item generation d'un support pack de soutenance.
\end{itemize}

La couverture, les scans Semgrep/Gitleaks/Trivy et le deploiement kind complet
ne doivent etre declares \texttt{TERMINE} que si leurs rapports d'execution
sont presents dans les artefacts correspondants.

\subsection{Elements prets mais dependants de l'environnement}

Les composants suivants sont implementes et scriptes, mais leur validation
complete depend de l'environnement d'execution :

\begin{itemize}
  \item declenchement Jenkins par webhook GitHub ;
  \item generation de SBOM par \texttt{Syft} ;
  \item signature et verification d'images par \texttt{Cosign} ;
  \item promotion par digest sans reconstruction ;
  \item deploiement verifie en mode \texttt{execute} ;
  \item activation de \texttt{metrics-server} ;
  \item observation effective des \textit{HorizontalPodAutoscaler} ;
  \item installation de \texttt{Kyverno} ;
  \item passage progressif des policies de \texttt{Audit} vers \texttt{Enforce}.
\end{itemize}

Ces points ne doivent pas etre presentes comme executes si les artefacts de
preuve correspondants ne sont pas disponibles. Ils doivent etre presentes comme
prets a etre rejoues, avec dependances explicites.

\subsection{Preuves a conserver}

Les preuves finales attendues sont les suivantes :

\begin{itemize}
  \item \texttt{artifacts/jenkins/github-webhook-validation.md} ;
  \item \texttt{artifacts/release/supply-chain-execute-summary.md} ;
  \item \texttt{artifacts/release/sign-summary.txt} ;
  \item \texttt{artifacts/release/verify-summary.txt} ;
  \item \texttt{artifacts/release/promotion-by-digest-summary.txt} ;
  \item \texttt{artifacts/release/promotion-digests.txt} ;
  \item \texttt{artifacts/sbom/sbom-index.txt} ;
  \item \texttt{artifacts/validation/cluster-security-addons.md} ;
  \item \texttt{artifacts/final/final-validation-summary.md} ;
  \item \texttt{artifacts/support-pack/*.tar.gz}.
\end{itemize}

\subsection{Formulation recommandee pour la soutenance}

La formulation recommandee est la suivante :

\begin{quote}
Le socle DevSecOps, Kubernetes et demonstration \texttt{demo} est operationnel
au niveau code, manifests et scripts. La preuve runtime complete reste
dependante de l'environnement : Docker accessible, cluster \texttt{kind},
registre OCI, cles Cosign, Syft, images signees, \texttt{metrics-server} et
\texttt{Kyverno}. Aucun resultat Jenkins, Cosign ou Kubernetes n'est declare
valide sans artefact d'execution.
\end{quote}
